\documentclass[12pt,compress,aspectratio=169]{beamer}
\input{../../mybeamer}

\title{Class 11: Rigid Body Rotational Motion, Part 3}
\subtitle{AP Physics C--Mechanics}
\input{../../me}
\input{../../mycommands}


\begin{document}

\begin{frame}
  \maketitle
\end{frame}



\section{Review}

%\begin{frame}{Curvilinear vs.\ Rectilinear Motion}
%  Kinematic quantities for rectilinear (translational) vs.\ curvilinear
%  (circular) motion are related:
%
%  \vspace{-.25in}{\large
%    \begin{align*}
%      \vec r &\quad\rightarrow\quad \theta \\
%      \vec v &\quad\rightarrow\quad \omega \\
%      \vec a &\quad\rightarrow\quad \alpha
%    \end{align*}
%  }
%
%  Dynamics:
%  
%  \vspace{-.25in}{\large
%    \begin{align*}
%      m &\quad\rightarrow\quad I\\
%      \vec F &\quad\rightarrow\quad\tau\\
%      \vec p=m\vec v_\text{cm} &\quad\rightarrow\quad L=I\omega
%    \end{align*}
%  }
%\end{frame}
%
%
%
%\begin{frame}{Laws of Motion}
%  The laws of motion are also related between translational and rotational
%  motion:
%  
%  \vspace{-.2in}{\large
%    \begin{align*}
%      \vec F_\text{net}=\diff{\vec p_\text{cm}}t &\quad\rightarrow\quad
%      \tau_\text{net}=\diff Lt \\
%      \vec F_\text{net}=m \vec a_\text{cm} &\quad\rightarrow\quad
%      \tau_\text{net}=I\alpha
%    \end{align*}
%  }
%\end{frame}




\begin{frame}{Translational Work}

  Translational work is done when a \emph{force} displaces an object by $\dl x$:

  \eq{-.1in}{
    W_\text{trans}=\int_{x_1}^{x_2}\vec F\cdot\dl\vec x
  }

  The net translational work done on an object (i.e.\ work done by the net
  force) is equal to the change in translational kinetic energy:

  \eq{-.15in}{
    W_\text{net,trans}=\Delta K_\text{trans}
  }

  \vspace{-.1in}where

  \eq{-.2in}{
    K_\text{trans}=\frac12mv_\text{cm}^2
  }

  It does not matter what $\vec F_\text{net}$ consists of. If $\vec F_\text{net}$
  does any work, it will \emph{always} be equal to the change in translational
  kinetic energy.
%  Mechanical Work:
%
%  \eq{-.1in}{
%    W_\text{trans}=\int_{x_1}^{x_2}\vec F\dl x
%    \quad\longrightarrow\quad
%    W_\text{rot}=\int_{\theta_1}^{\theta_2}\tau\dl\theta
%  }
%
%  Work-energy theorem:
%
%  \eq{-.1in}{
%    W_\text{trans}=\Delta K_\text{trans}
%    \quad\longrightarrow\quad
%    W_\text{rot}=\Delta K_\text{rot}
%  }
%
%  Kinetic Energy:
%  
%  \eq{-.1in}{
%    K_\text{trans}=\frac12mv_\text{cm}^2
%    \quad\longrightarrow\quad
%    K_\text{rot}=\frac12I\omega^2
%  }
\end{frame}



\begin{frame}{Rotational Work}

  Likewise, rotational work is done when a \emph{torque} displaces an object by
  $\dl\theta$ along the direction of rotation:

  \eq{-.2in}{
    W_\text{rot}=\int_{\theta_0}^{\theta_1}\tau\dl\theta
  }

  The net rotational work done on an object (i.e.\ work done by the net torque)
  is equal to the change in rotational kinetic energy:

  \eq{-.15in}{
    W_\text{net,rot}=\Delta K_\text{rot}
  }

  \vspace{-.1in}where

  \eq{-.2in}{
    K_\text{rot}=\frac12I\omega^2
  }

  \vspace{-.1in}It does not matter what $\tau_\text{net}$ consists of; as long as
  $\tau_\text{net}$ does work, it will \emph{always} be equal to the change in
  rotational kinetic energy.
\end{frame}



\begin{frame}{Kinetic Energy of a Rotating System}
  The total kinetic energy of a rotating system is the sum of its translational
  and rotational kinetic energies at its centre of mass:

  \eq{-.1in}{
    \boxed{
      K_\text{tot}
      =K_\text{trans}+K_\text{rot}
      =\frac12mv_\text{cm}^2+\frac12I_\text{cm}\omega^2}
  }
  
  In this case, $I_\text{cm}$ is calculated at the centre of mass. For simple
  problems, we only need to compute rotational kinetic energy at the pivot:
  
  \eq{-.1in}{
    \boxed{K_\text{rot}=\frac12I_\text{pivot}\omega^2}
  }
  
  In this case, the $I_\text{pivot}$ is calculated at the pivot.
  \textbf{IMPORTANT:} $I_\text{cm}\neq I_\text{pivot}$
\end{frame}



\begin{frame}{Conservative Forces}
  Forces that are related to a potential energy are called
  \textbf{conservative forces}. They include:
  \begin{itemize}
  \item Gravitational force $\vec F_g$
  \item Spring force $\vec F_s$
  \item Electrostatic force $\vec F_q$
  \item Magnetic force $\vec F_m$
  \item Nuclear forces
  \end{itemize}
\end{frame}




\begin{frame}{Conservative Forces}
  All conservative forces share these common properties:
  \begin{enumerate}
  \item The work done by these forces relate to a change of a potential energy.
    This is often called the \textbf{work-potential energy theorem}:

    \eq{-.1in}{
      \boxed{
        W_\text{conservative}=-\Delta U
      }
    }
    \begin{itemize}
    \item\vspace{-.17in}Positive work decreases this related potential energy
    \item Negative work increases this related potential energy
    \end{itemize}
  \item $+$ work done by a conservative force converts potential energy
    into kinetic energy
  \item $-$ work done by a conservative force converts kinetic energy
    into potential energy
  \item The work done by a conservative force is \emph{path independent}
    %(depends only on end points)
  \end{enumerate}
\end{frame}



\begin{frame}{Conservative Forces}
  By the fundamental theorem of calculus, any conservative forces $\vec F$
  must be the negative gradient of the potential energies:

  \eq{-.1in}{
    \boxed{
      \vec F=-\nabla U=
      -\diffp Ux\hat\imath-\diffp Uy\hat\jmath-\diffp Uz\hat k
    }
  }

  In one-dimension (which is likely to be what you may encounter in AP Physics
  C):

  \eq{-.1in}{
    \boxed{
      F=-\diff Ux
    }
  }

  The direction of a conservative force \emph{always} decreases the potential
  energy. (Pay attention to the negative sign.)
\end{frame}



\begin{frame}{Examples of Non-Conservative Force}
  The majority of forces are \textbf{non-conservative}. The common forces
  discussed in the previous class are generally non-conservative:
  \begin{itemize}
  \item Applied force
  \item Tension force
  \item Normal force
  \item Friction
  \item Drag (fluid resistance)
  \end{itemize}
  The work-energy theorem still applies for non-conservative forces
\end{frame}



\begin{frame}{Work by Non-Conservative Forces}
  The work done by non-conservative forces differs from conservative forces in
  that:
  \begin{itemize}
  \item There is \textbf{no related potential energies}: the work done by a
    non-conservative force transforms energy
    \begin{itemize}
    \item From kinetic energy of one object to another
    \item From one form of kinetic energy to another
    \item From kinetic energy into thermal energy
    \end{itemize}
  \item The work is generally \textbf{path dependent}
  \end{itemize}
\end{frame}




\section{Rolling on Inclined Surface}

\begin{frame}{Rolling on an Inclined Surface}
  We now re-examine the example problem for a rigid \& smooth sphere of radius
  $R$ rolling down a ramp of angle  $\phi$ without slippage, this time using
  the conservation of energy.

  \vspace{.2in}
  \begin{columns}
    \column{.35\textwidth}
    \input{../Class10-rotMotion-part2/roll-down-ramp}

    \column{.65\textwidth}
    Three forces act on the sphere as it rolls down the ramp
    \begin{itemize}
    \item The weight ($\vec F_g$) of the sphere acts at the CM
    \item The normal force ($\vec F_n$) acts at the point of contact
    \item The static friction ($\vec f_s$) act at the point of contact
    \end{itemize}
    Only static friction generates a torque about the CM in the clockwise
    direction
  \end{columns}
\end{frame}



\begin{frame}{Rolling on an Inclined Surface}
  The centre of mass of the sphere will translate down the ramp, along the
  $\iii$ direction. The sphere will also rotate in the clockwise direction.

  \vspace{.2in}
  \begin{columns}
    \column{.35\textwidth}
    \input{../Class10-rotMotion-part2/roll-down-ramp}

    \column{.65\textwidth}
    The work done by the three forces acting on the sphere:
    \begin{itemize}
    \item Weight ($\vec F_g$) does $+$ translational work, but no rotational
      work
    \item Normal force ($\vec F_n$) does no translational work, and no
      rotational work
    \item Static friction ($\vec f_s$) does $(-)$ translational work, and $(+)$
      rotational work
    \end{itemize}
    Although static friction is non-conservative, the work done stays inside
    the system
  \end{columns}    
\end{frame}



\begin{frame}{Rolling on an Inclined Surface}
  In this case, the system consists of the the sphere and Earth, much like the
  (much simpler) system of a mass sliding down a ramp.

  \vspace{.2in}
  \begin{columns}
    \column{.35\textwidth}
    \input{../Class10-rotMotion-part2/roll-down-ramp}

    \column{.65\textwidth}
    The total energy of the system therefore includes:
    \begin{itemize}
    \item The translational kinetic energy of the sphere
    \item The rotational kinetic energy of the sphere
    \item The gravitational potential energy stored between the sphere \& Earth
    \end{itemize}
    The system is an isolated system, because there is no external work done
    \begin{itemize}
    \item Although friction is non-conservative, it is \emph{internal} to the
      system
    \end{itemize}
  \end{columns}
\end{frame}



\begin{frame}{Rolling on an Inclined Surface}
  \begin{columns}
    \column{.35\textwidth}
    \input{../Class10-rotMotion-part2/roll-down-ramp}

    \column{.65\textwidth}
    As the sphere rolls downward without slipping, the total mechanical
    energy is conserved, i.e.:

    \eq{-.1in}{
      \boxed{
        K_\text{trans} + K_\text{rot} + U_g = \text{constant}
      }
    }
    \begin{itemize}
    \item\vspace{-.1in}The positive translational work done by gravity converts
      $U_g$ into $K_\text{trans}$
    \item The work done by static friction converts $K_\text{trans}$ to
      $K_\text{rot}$
    \item All the energy stays inside the system, therefore the system is an
      isolated system
    \end{itemize}
  \end{columns}
\end{frame} 
\end{document}

